%voorwoord en dankwoord
Voor u ligt de thesis "Ontwikkeling van een smartphone SDK voor het obfusceren van locatie data". Deze thesis is geschreven om te voldoen aan de afstudeereisen van de Master in de
industriële wetenschappen elektronica-ICT:
smart applications aan de Katholieke Universiteit van Leuven.

Tijdens het schrijven van deze thesis heb ik veel geleerd over zowel de technische aspecten van locatie-obfuscatie als over het zelfstandig uitvoeren van onderzoek. Het was een uitdagend proces waarin ik mijn kennis over softwareontwikkeling, privacy en mobiele applicaties heb kunnen verdiepen. Zeker de diepgang in de werking van het Android besturingssysteem was een uiterst interessant onderwerp, verder waren de experimenten in deze thesis uniek doordat ik buiten kon gaan wandelen, fietsen,... en hieruit rechtstreeks enorm visuele resultaten kon bekomen. Dit zorgde voor directe validatie en motivatie.

Graag maak ik van deze gelegenheid gebruik om enkele mensen te bedanken die een belangrijke rol hebben gespeeld bij het tot stand komen van deze thesis.

In de eerste plaats wil ik mijn promotor, prof. dr. Vincent Naessens, bedanken voor het vertrouwen in dit onderzoek en voor het aanreiken van dit boeiende onderwerp. Mijn oprechte dank gaat ook uit naar mijn masterproefbegeleider, dr. ing. Michiel Willocx, voor de deskundige begeleiding, waardevolle feedback en voortdurende motivatie gedurende het hele traject. Daarnaast wil ik de medewerkers van Distrinet bedanken voor hun steun en de aangename gesprekken doorheen het jaar.

Verder ben ik mijn vrienden en familie, Bran, Cas, Melissa, Robin, Warre, 2x Lotte, Angelo, Elise, Sofie en vele anderen dankbaar voor hun vriendschap, aanmoediging en de mooie jaren samen in het hoger onderwijs.

Tot slot wil ik iedereen bedanken die mij op welke manier dan ook heeft gesteund tijdens mijn masterjaar. Jullie steun heeft het verschil gemaakt.

Ik hoop dat het onderwerp u even hard kan boeien als mij.

Milan Schollier

Gent, \today